\chapter{The free Maxwell comparison theory}\label{ch03:maxwell}

The reference theory for short distances consists of independent Maxwell
fields, one for each direction in the Lie algebra. Its role is more precise
than a formal approximation to the action. It supplies finite-dimensional
spaces of local fields, correlation functionals that distinguish those
fields, and a normalization for the stress tensor. The interacting
construction must eventually reproduce these data in an appropriate
short-distance comparison. No interacting limit is assumed in this chapter.

We fix the positive inner product on the real Lie algebra $\gLie$ and write
$d_{\gLie}=\dim_{\RR}\gLie$. Use the orthonormal basis of Section~\ref{ch01:gauge},
$(e_\alpha)_{\alpha=1}^{d_{\gLie}}$. Greek letters $\mu,\nu,\rho,\sigma$
from the middle of the alphabet denote coordinate indices in
$\{1,2,3,4\}$; $\alpha,\beta$ denote indices in this basis. Time is the fourth
coordinate. The superscript $\mathrm{Max}$ always denotes the free
reference theory. We first define its correlations at distinct points;
this requires no multiplication of distributions at the same point.

\section{Curvature covariance}\label{ch03:covariance}

Let
\begin{equation}\label{ch03:green}
 \Gamma(x)=\frac{1}{4\pi^2|x|^2},\qquad x\ne0.
\end{equation}
Its locally integrable extension is the fundamental solution of
$-\Delta$ in four dimensions. A convenient way to specify the curvature
covariance is to start from independent Gaussian vector-potential
components with covariance
$\delta_{\alpha\beta}\delta_{\mu\nu}\Gamma(x-y)$, and take a curl in
each argument. The vector potential in this description is a device for
computing the covariance; it is not one of the gauge-invariant local
observables used in the construction.

Write $F^\alpha_{\mu\nu}=-F^\alpha_{\nu\mu}$ for the curvature
components in this basis. The covariance is the product of the inner
product on $\gLie$ and a spacetime kernel $C$: pairing the two fields
with $u,v\in\gLie$ gives $\langle u,v\rangle_{\gLie}C_{\mu\nu,\rho\sigma}(x-y)$.
In coordinates,
\begin{equation}\label{ch03:curvature}
 S^{\mathrm{Max}}_2(F^\alpha_{\mu\nu},F^\beta_{\rho\sigma};x,y)
 =\delta_{\alpha\beta}C_{\mu\nu,\rho\sigma}(x-y),
\end{equation}
where, with $h_{\mu\nu}=\partial_\mu\partial_\nu\Gamma$,
\begin{align}
 C_{\mu\nu,\rho\sigma}
 &=-\delta_{\nu\sigma}h_{\mu\rho}
   +\delta_{\nu\rho}h_{\mu\sigma}
   +\delta_{\mu\sigma}h_{\nu\rho}
   -\delta_{\mu\rho}h_{\nu\sigma},\label{ch03:kernel}\\
 h_{\mu\nu}(z)
 &=\frac{8z_\mu z_\nu-2\delta_{\mu\nu}|z|^2}
          {4\pi^2|z|^6}.\label{ch03:hessian}
\end{align}
The minus signs in the first and last terms arise because a derivative
in $y$ is minus a derivative in $z=x-y$. This convention matters both
for reflection positivity and for the stress normalization.

The covariance is translation invariant and transforms as a two-form in
each pair of indices under $O(4)$. At nonzero separation it satisfies
\begin{equation}\label{ch03:maxwell-equations}
 \partial_\lambda C_{\mu\nu,\rho\sigma}
 +\partial_\mu C_{\nu\lambda,\rho\sigma}
 +\partial_\nu C_{\lambda\mu,\rho\sigma}=0,
 \qquad
 \sum_\mu\partial_\mu C_{\mu\nu,\rho\sigma}=0.
\end{equation}
The first identity is the Bianchi identity; the second is the source-free
Maxwell equation. They hold off the origin. A distributional derivative
can have a contact term at the origin, so these identities must not be
used to infer a choice of contact extension.

For a multi-index $a=(a_1,\ldots,a_4)$, put
$|a|=a_1+\cdots+a_4$ and
$\partial^a=\partial_1^{a_1}\cdots\partial_4^{a_4}$.
Derivatives of the fields have covariance
\begin{equation}\label{ch03:jet-covariance}
 S^{\mathrm{Max}}_2(\partial^aF^\alpha_{\mu\nu},
                   \partial^cF^\beta_{\rho\sigma};x,y)
 =\delta_{\alpha\beta}(-1)^{|c|}
       \partial^{a+c}C_{\mu\nu,\rho\sigma}(x-y).
\end{equation}
In particular, for each fixed derivative order, this kernel is smooth
away from zero and homogeneous of degree $-4-|a|-|c|$. Thus its size on
$|x-y|\ge r$, $0<r\le1$, has a polynomial bound in $r^{-1}$.
This elementary estimate is the free counterpart of the much harder
uniform bounds required from the lattice construction.

\section{Local polynomials and Wick correlations}\label{ch03:wick}

A local polynomial is a polynomial in finitely many curvature derivatives
at one point. Its scaling dimension is computed by assigning degree $2$
to $F$ and degree $1$ to each derivative. For example,
$F^\alpha_{\mu\nu}F^\beta_{\rho\sigma}$ has degree $4$, whereas
$\partial_\lambda F^\alpha_{\mu\nu}$ has degree $3$.
Scaling dimension is an algebraic grading of the reference model; it is
not an assertion about anomalous dimensions of the interacting fields.

The polynomial attached to one insertion is Wick ordered. We use the
usual Gaussian pairing rule, whose operator formulation is due to
Wick~\cite{Wick1950}. Here it can be taken as the definition of the
separated correlation functions. Expand each polynomial into monomials,
regard each factor as a distinguishable occurrence, and sum over pairings
of all occurrences subject to the condition that paired occurrences
belong to different insertions. For each pair use
\eqref{ch03:jet-covariance}. The empty pairing has value one. The
resulting finite sum is denoted by
\[
 S_n^{\mathrm{Max}}(P_1,\ldots,P_n;x_1,\ldots,x_n),
 \qquad x_i\ne x_j\quad(i\ne j).
\]
There is no covariance at zero separation in this definition.

For instance, if $X_1,X_2,Y_1,Y_2$ are linear curvature derivatives,
\begin{align*}
 S_2^{\mathrm{Max}}({:}X_1X_2{:},{:}Y_1Y_2{:};x,y)
 &=C(X_1,Y_1;x-y)C(X_2,Y_2;x-y)\\
 &\quad+C(X_1,Y_2;x-y)C(X_2,Y_1;x-y).
\end{align*}
Here $C(X,Y;z)$ denotes the bilinear covariance obtained from
\eqref{ch03:jet-covariance}. By contrast, in a word containing several
insertions in the same time half-space, contractions between those
\emph{different} insertions are included. Wick ordering removes
self-contractions at each insertion, not all contractions inside a word.

The pairing rule gives complex multilinearity, symmetry under simultaneous
permutation of labels and positions, reality, and the unit identity
\[
 S_{n+1}^{\mathrm{Max}}(1,P_1,\ldots,P_n;x_0,x_1,\ldots,x_n)
 =S_n^{\mathrm{Max}}(P_1,\ldots,P_n;x_1,\ldots,x_n).
\]
Differentiating an insertion differentiates its incident covariances.
The product rule therefore agrees with ordinary differentiation of the
correlation function in that insertion position.

If the labels are homogeneous of degrees $d_1,\ldots,d_n$, then
\begin{equation}\label{ch03:homogeneity}
 S_n^{\mathrm{Max}}(P_1,\ldots,P_n;\lambda x_1,\ldots,\lambda x_n)
 =\lambda^{-\sum_i d_i}
 S_n^{\mathrm{Max}}(P_1,\ldots,P_n;x_1,\ldots,x_n)
 \quad(\lambda>0).
\end{equation}
For general labels of bounded degree, decompose into finitely many
homogeneous pieces. Equations \eqref{ch03:jet-covariance} and
\eqref{ch03:homogeneity} give bounds on all derivatives on separated
configuration spaces. In particular, smearing against a smooth compactly
supported function whose support avoids collisions is well defined.
This statement is weaker than specifying the correlations on test functions whose supports meet collision diagonals.

\section{The on-shell polynomial space}\label{ch03:shell}

The Maxwell identities impose linear relations on the jet variables.
They should be imposed before discussing strict positivity or the
number of independent fields. At each integer $j\ge0$, let
$\mathcal S_j$ be the real vector space of homogeneous degree-$j$
polynomial two-forms on $\RR^4$, with values in $\gLie$, satisfying
\[
 dF=0,\qquad \delta F=0.
\]
Here $\delta F=0$ means
$\sum_\mu\partial_\mu F_{\mu\nu}=0$ for every $\nu$; the harmless
choice of sign for the codifferential does not affect its kernel.
These are finite-dimensional spaces. They encode the possible homogeneous
Taylor coefficients of a source-free Maxwell curvature. A Taylor
coefficient of order $j$ is normalized as $D^jF(0)/j!$.

The phrase \emph{on shell} will mean restriction to these solution jets.
It does not refer to restricting a point of Euclidean spacetime.
A linear coordinate on $\mathcal S_j$ has scaling dimension $j+2$. As in Definition~\ref{ch02:labels}, we also call
`weight' the scaling dimension of a homogeneous jet polynomial.
Consequently polynomials of weight at most $d$ use only the spaces
$\mathcal S_j$ with $j\le d-2$. Constants are retained at every degree.
The action on $\gLie$ is the adjoint action of $G$; invariant polynomials
are the gauge-invariant local observables of the free reference theory.

To compare this definition with the radial-potential labels $V_d$,
associate to a closed homogeneous curvature $F^{(j)}$ the radial
potential
\begin{equation}\label{ch03:radial}
 A^{(j+1)}_\mu(x)=\int_0^1 t x^\nu F^{(j)}_{\nu\mu}(tx)\,dt
 =\frac{1}{j+2}\sum_\nu x^\nu F^{(j)}_{\nu\mu}(x).
\end{equation}
There is a sum over $\nu$ in the integral as well. This is the usual
radial homotopy formula: $x^\mu A_\mu=0$ and $dA=F$ when $dF=0$.
The connection here is the \emph{linear} Maxwell potential; no
nonabelian commutator is inserted into this formula. Restricting a label
in $V_d$ to these radial potentials gives a map
\[
 \omega_d:V_d\longrightarrow
 \{
 \text{complex polynomials on Maxwell solution jets of weight at most }d
 \}.
\]
It preserves the unit, conjugation and the geometric actions. Define
\begin{equation}\label{ch03:free-space}
 \mathcal M_d^{\mathrm{Max}}=\operatorname{im}\omega_d.
\end{equation}
Equivalently, this is $V_d/\ker\omega_d$ as a vector space. The image
formulation is useful: it does not assume a separate surjectivity theorem
onto a larger invariant-polynomial space.

The covariance identities \eqref{ch03:maxwell-equations} make the Wick
correlations well defined on these on-shell polynomial images. Products
of a Maxwell equation with another jet polynomial vanish in every
separated Wick correlation. Conversely, proving that \emph{all}
correlation-null polynomials vanish on shell requires a nondegeneracy
argument; it does not follow from writing down the equations.

The free spaces form a filtration. The inclusion
$\mathcal M_d^{\mathrm{Max}}\to\mathcal M_e^{\mathrm{Max}}$ for
$d\le e$ retains the same polynomial, while differentiation raises
weight by one. For a simple nonabelian Lie algebra there is no nonzero
invariant linear form on $\gLie$. Therefore every gauge-invariant field
of degree less than four in this polynomial image is a constant.
This low-degree statement follows from the weights: a nonconstant
monomial below degree four would have just one curvature factor or one
of its first derivatives.

\section{Reflection and the meaning of nondegeneracy}\label{ch03:reflection}

Let $\vartheta(\boldsymbol x,t)=(\boldsymbol x,-t)$.
Reflection acts on a two-form by pullback, so a component with one time
index changes sign and a purely spatial component does not. On a local
polynomial, $\theta$ combines this action with complex conjugation.
On a word it also reverses the order and reflects every insertion point.
This is the antilinear reflection used in the Osterwalder--Schrader
form, not merely conjugation of scalar coefficients.

For words $w_i$ supported at distinct points in the positive half-space,
the free field has reflection positivity:
\begin{equation}\label{ch03:rp}
 \sum_{i,j}\overline{c_i}c_j
 S^{\mathrm{Max}}(\theta w_i\,w_j)\ge0.
\end{equation}
The same inequality holds for smooth joint test functions supported on separated
positive-time configurations. This is the standard Gaussian
reflection-positive construction underlying
Osterwalder--Schrader reconstruction~\cite{OsterwalderSchrader1973,OsterwalderSchrader1975},
applied to the curvature covariance rather than to an unconstrained
vector-potential observable. Its hypotheses can be checked at the
one-particle level: spatial Fourier transformation of
\eqref{ch03:green} gives
\[
 \frac{e^{-|\boldsymbol p|(t+s)}}{2|\boldsymbol p|}
 \qquad(t,s>0),
\]
and taking the two curls with the reflected two-form indices gives the
positive transverse-polarization Gram form. Derivative jets act on the
same factors. The Gaussian Fock construction then gives
\eqref{ch03:rp}; contractions within each positive-time word are included
in its Fock vector. On compact separated supports the kernels are smooth,
so the corresponding form on joint test functions follows by integration. We use this
standard argument by citation rather than repeating the Fock-space
construction.

A simple sign check is informative. At separation $z=re_4$, the six
unordered curvature components have diagonal covariance. Three purely
spatial components have value $1/(\pi^2r^4)$; the three components
$F_{i4}$ have value $-1/(\pi^2r^4)$. Reflection changes the sign of
$F_{i4}$ in the first slot. The reflected one-jet form is therefore
positive in all six directions. Without the tensorial reflection, the
same displayed covariance would give the wrong positivity statement.

Positive semidefiniteness of \eqref{ch03:rp} alone does not identify its
null space. For later rank arguments one needs the stronger statement
that a nonzero on-shell polynomial has a strictly positive reflected
two-point pairing at two suitable points. Derivatives of different
orders may mix in that pairing. Checking one homogeneous order at a
time does not establish the statement for their sum.

\section{Strict positivity of finite jet families}\label{ch03:strict}

We now identify the null relations in the free polynomial theory.
Recall from Section~\ref{ch03:shell} that $\mathcal S_j$ is the space
of homogeneous degree-$j$ Maxwell curvature solutions. Write a member as
\[
 P_{ab}(v)=\sum_{i_1,\ldots,i_j}
 P_{ab;i_1\cdots i_j}v_{i_1}\cdots v_{i_j}.
\]
The derivative-index coefficients are symmetric. On this coefficient
space use the Frobenius inner product, summing both orders of the
antisymmetric field indices, all ordered derivative indices and components in the basis of Section~\ref{ch01:gauge}.
In particular, no factor $1/2$ is inserted for the field indices.
The linear jet probe represented by $P$ is
\begin{equation}\label{ch03:probe}
 \ell_P(F)=\frac1{j!}\sum_{a,b,i_1,\ldots,i_j,\alpha}
 P^\alpha_{ab;i_1\cdots i_j}
 \partial_{i_1}\cdots\partial_{i_j}F^\alpha_{ab}.
\end{equation}
The factor $1/j!$ matches our Taylor convention. Every linear functional
on the order-$j$ solution jets has this form.

Put $I(z)=z/|z|^2$ and
$R_z=1-2zz^{\mathsf T}/|z|^2$. For a two-form, $R_z^{(2)}$ means
applying $R_z$ to both form indices. The full reflection on a
polynomial solution is
$(R_z\cdot P)(v)=R_z^{(2)}P(R_zv)$: the derivative indices are
reflected as well. For real probes define
\[
 B_z(P,Q)=S_2^{\mathrm{Max}}(\ell_P,\ell_{R_z\cdot Q};z,0).
\]
Extend this bilinearly across different derivative orders. Covariance and
exchange of the two insertions show that $B_z$ is symmetric.

\begin{proposition}[Finite-jet nondegeneracy]\label{ch03:jet-positive}
For every $z\ne0$ and integer $m\ge1$, the form $B_z$ is positive
definite on $\bigoplus_{j<m}\mathcal S_j$.
\end{proposition}

The scalar differentiation identity needed in the proof is the harmonic
case of Hobson's differentiation formula~\cite{Hobson1892}. In our
normalization, if $P\in\mathcal S_j$, it says, componentwise,
\begin{equation}\label{ch03:hobson}
 \frac1{j!}P(\partial)\Gamma(y)
   =\frac{(-2)^j}{4\pi^2}|y|^{-2j-2}P(y),\qquad y\ne0.
\end{equation}
Here $P(\partial)$ means the ordered coefficient contraction without
the factor $1/j!$. Maxwell solutions are harmonic, as follows by
applying a divergence to the Bianchi identity. Thus the harmonic
hypothesis of the scalar formula holds in each component. We use that
classical differentiation result; the following finite-dimensional
argument explains its application to the mixed jet orders needed here.

\begin{proof}
Let $\mathcal R(P)(y)$ be the curvature response obtained by contracting
$\partial_y^j C(y)/j!$ with the coefficients of $P$. Taking the two
curls in \eqref{ch03:hobson} gives
\begin{equation}\label{ch03:kelvin}
 \mathcal R(P)(y)=c_j|y|^{-4}R_y^{(2)}P(I(y)),
 \qquad c_j=\frac{(j+1)(j+2)(-2)^j}{\pi^2}.
\end{equation}
For clarity, the identities determining the coefficient are as follows.
Set $V_i(y)=\sum_a y_aP_{ai}(y)$ and
$b_j=(-2)^j/(4\pi^2)$. The divergence of the differentiated Green
potential is
$D_i=-2(j+1)b_j|y|^{-2j-4}V_i$.
The Maxwell equations and homogeneity give
$\partial_iV_m-\partial_mV_i=(j+2)P_{im}$.
The response is $2(\partial_mD_i-\partial_iD_m)$, so its coefficient
is $4(j+1)(j+2)b_j=c_j$ and its two-form factor is
\[
 P_{im}(y)-2|y|^{-2}(y_iV_m(y)-y_mV_i(y))
   =(R_y^{(2)}P(y))_{im}.
\]
This verifies \eqref{ch03:kelvin}, including its sign, off the origin.

Fix the separation $z$. For $P\in\mathcal S_j$ define a finite tower
$\Sigma_zP$ by
\begin{equation}\label{ch03:triangular}
 (\Sigma_zP)_k
   =\frac{c_j}{c_k}\operatorname{Taylor}_k[P(\,\cdot-I(z))],
 \qquad 0\le k\le j.
\end{equation}
The operator $\operatorname{Taylor}_k$ extracts the homogeneous term of
degree $k$ at zero. Translation preserves the Maxwell equations, and
each homogeneous term separately satisfies them. Thus all components in
\eqref{ch03:triangular} lie in their declared solution spaces.
The degree-$j$ component is $P$ itself.

By the finite Taylor formula, the response of this tower is
\[
 \mathcal R(\Sigma_zP)(y)
   =c_j|y|^{-4}R_y^{(2)}P(I(y)-I(z)).
\]
Since $DI(z)=|z|^{-2}R_z$, its Taylor coefficients at $y=z$ vanish
through order $j-1$, while its order-$j$ coefficient is
$c_j|z|^{-2j-4}(R_z\cdot P)$.
For $Q\in\mathcal S_k$, differentiation in the second insertion gives
\[
 B_z(\Sigma_zP,Q)
 =(-1)^k\big\langle R_z\cdot Q,
       \operatorname{Taylor}_k[\mathcal R(\Sigma_zP)](z)
                 \big\rangle_F.
\]
It follows that this expression vanishes if $k<j$, and at $k=j$ equals
$(-1)^jc_j|z|^{-2j-4}\langle P,Q\rangle_F$.
Transforming the second argument introduces only lower orders in addition
to its unchanged top order. Symmetry therefore gives, on the full tower,
\begin{equation}\label{ch03:diagonal-gram}
 B_z(\Sigma_zP,\Sigma_zQ)
 =\sum_{j<m}\frac{(j+1)(j+2)2^j}{\pi^2|z|^{2j+4}}
                  \langle P_j,Q_j\rangle_F.
\end{equation}
Every coefficient is positive. Finally $\Sigma_z$ is invertible:
it is the identity plus a strictly degree-lowering map, hence its inverse
is a finite geometric sum. This proves positivity on the whole tower,
including mixtures of derivative orders.
\end{proof}

\begin{corollary}[Polynomial positivity and finite separating families]
\label{ch03:detectors}
Every nonzero element of $\mathcal M_d^{\mathrm{Max}}$ has strictly
positive reflected two-point pairing at a fixed separated pair.
For any basis $(P_1,\ldots,P_s)$ of this space and fixed distinct points
$x,y$, the linear map
\[
 P\longmapsto
 \big(S_2^{\mathrm{Max}}(P,P_j;x,y)\big)_{j=1}^s
\]
is injective. Consequently a label in $V_d$ vanishes in every separated
correlation if and only if its on-shell polynomial $\omega_d$ vanishes.
\end{corollary}

\begin{proof}
This is the standard symmetric-Fock-space consequence of the one-particle
Gram form and Wick's rule~\cite{Wick1950}. Explicitly, choose coordinates
$X_1,\ldots,X_M$ diagonalizing the positive form on a tower containing
all variables of the polynomial, with positive squared lengths $b_i$.
The reflected pairing of Wick monomials is
\[
 \langle X^a,X^c\rangle_{\rm ref}
   =\mathbf1_{a=c}\,a!\prod_i b_i^{a_i}.
\]
Thus the squared norm of a polynomial is the sum of the squared absolute
values of its coefficients with these positive weights; the constant
coefficient has weight one. Restricting to the gauge-invariant image
preserves positivity. If all basis-column pairings vanish, the pairing
with the reflected conjugate of the polynomial vanishes too, forcing
that polynomial to be zero. This also identifies the correlation-null
labels. On any fixed finite-dimensional degree space, the map of correlation pairings
has a positive minimum norm on the unit sphere. Small perturbations of
its matrix remain injective; applying this fact to interacting
correlations still requires a quantitative comparison estimate.
\end{proof}


\section{Quadratic fields and normalization checks}\label{ch03:quadratic}

We return to the free curvature covariance
\eqref{ch03:curvature}; in the basis of Section~\ref{ch01:gauge}, each component has the same
covariance and orthogonal components have zero cross-covariance. Define
\begin{align}
 Q_{\mu\nu}
 &=\sum_{\rho=1}^4
       {:}\langle F_{\mu\rho},F_{\nu\rho}\rangle_{\gLie}{:},
 & Q_{\mathrm{tot}}&=\sum_\mu Q_{\mu\mu},\label{ch03:quadratics}\\
 \mathcal E^{\mathrm{Max}}&=\tfrac14 Q_{\mathrm{tot}},
 &T^{\mathrm{Max}}_{\mu\nu}
   &=-Q_{\mu\nu}+\tfrac14\delta_{\mu\nu}Q_{\mathrm{tot}}.
 \label{ch03:stress}
\end{align}
The sign of $T^{\mathrm{Max}}$ is chosen to match the shell-flux
convention in the supplied construction. A complete stress-normalization
argument must check that convention against the translation Ward identity;
conservation alone does not determine the sign.

The tensor in \eqref{ch03:stress} is symmetric and traceless. Its
divergence vanishes as an on-shell Wick polynomial. These are the
ordinary algebraic Maxwell stress identities, obtained from the Bianchi
and source-free equations; the same identities hold in separated
correlations. The scalar $\mathcal E^{\mathrm{Max}}$ is nevertheless
nonzero. It is the curvature-square operator, not the trace of the free
stress tensor.

With the present normalization, Wick's rule gives
\begin{equation}\label{ch03:scalar-pair}
 S_2^{\mathrm{Max}}(\mathcal E^{\mathrm{Max}},
                    \mathcal E^{\mathrm{Max}};x,y)
 =\frac{3d_{\gLie}}{\pi^4|x-y|^8},\qquad x\ne y.
\end{equation}
For checking constants, at separation $re_4$ one has
$\mathcal E^{\mathrm{Max}}=\tfrac12\sum_{\mu<\nu}{:}\|F_{\mu\nu}\|_{\gLie}^2{:}$. The two cross-pairings of each square,
the six unordered components, and the prefactor $1/4$ give
$\tfrac14\cdot2\cdot6d_{\gLie}=3d_{\gLie}$.
Orthogonal covariance makes the answer independent of the direction of
$x-y$. The calculation uses only separated Wick contractions.

This normalization will be needed when comparing a renormalized scalar
field with the free curvature square. It does not itself identify the
interacting trace, its coefficient, or its logarithmic asymptotic.
Those conclusions require the interacting stress Ward identity and
quantitative estimates on moving comparison fields.

\section{The degree-four conserved tensor}\label{ch03:classification}

The finite-dimensional classification needed for the stress tensor is
small enough to state explicitly. In this section \emph{centered} means
that the constant Wick coefficient is zero. By the Wick prescription it
is also the one-point expectation. We consider real symmetric tensor
fields $J_{ab}$ of degree at most four, invariant under the adjoint action
and transforming as a rank-two tensor under the signed coordinate
permutations $H_4$. Conservation is an identity in the on-shell polynomial
space, or equivalently in every separated free correlation by
Corollary~\ref{ch03:detectors}.

\begin{proposition}[Conserved tensor line]\label{ch03:stress-line}
The space of such centered tensors has dimension three. Its conserved
subspace is the line $\RR T^{\mathrm{Max}}$, with
$T^{\mathrm{Max}}$ as in \eqref{ch03:stress}.
\end{proposition}

\begin{proof}
Simplicity of $\gLie$ excludes invariant linear forms on $\gLie$ and makes
an invariant symmetric bilinear form on $\gLie$ a multiple of the fixed
metric. Through degree four a centered invariant polynomial therefore
contains two undifferentiated curvature factors. Indeed, a derivative on
one of two factors would raise the degree to at least five, while a
single differentiated factor still has no invariant linear functional on $\gLie$.

Recall the quadratic fields $Q_{ab}$ and $Q_{\mathrm{tot}}$ from
\eqref{ch03:quadratics}. A basis of the tensor candidates is
\begin{align*}
 Z^{(1)}_{ab}&=\delta_{ab}Q_{aa},\\
 Z^{(2)}_{ab}&=\delta_{ab}(Q_{\mathrm{tot}}-2Q_{aa}),\\
 Z^{(3)}_{ab}&=(1-\delta_{ab})Q_{ab}.
\end{align*}
There is no sum over $a$ on the right sides of the first two formulas.
To see that these span, consider sign reversal of each coordinate.
A diagonal component can contain only squares of curvature components.
Permutations fixing its distinguished axis divide them into two classes:
those containing that axis, and those containing neither instance of it.
An off-diagonal component with indices $a,b$ must be odd in axes $a,b$
and even in the others. Its possible contractions are
${:}\langle F_{ac},F_{bc}\rangle_{\gLie}{:}$ with $c\notin\{a,b\}$;
permutations identify their coefficients. This gives exactly the three
candidates displayed above.

For an explicit rank check, let
$Q_{ij,ij}={:}\|F_{ij}\|_{\gLie}^2{:}$.
Evaluate the first component of the divergence of each $Z^{(l)}$ at
$e_1$, against $Q_{12,12}(0)$ and $Q_{23,23}(0)$. The two rows of
these evaluations, after division by $d_{\gLie}/(4\pi^2)^2$, are
\begin{equation}\label{ch03:divergence-matrix}
 \begin{pmatrix}-256&0&128\\0&-512&-128\end{pmatrix}.
\end{equation}
These numbers come directly from the curvature covariance
\eqref{ch03:kernel} and the two cross-pairings in a quadratic Wick
correlation. For example,
$S_2^{\mathrm{Max}}(Q_{11},Q_{12,12};te_1,0)
=2d_{\gLie}/(\pi^4t^8)$; its derivative at $t=1$ gives the first
entry after the stated division. The first two columns are independent.
Also $Z^{(1)}_{12}=Z^{(2)}_{12}=0$, whereas the correlation of
$Z^{(3)}_{12}(e_1)$ with
${:}\langle F_{13},F_{23}\rangle_{\gLie}{:}(0)$ equals
$-16d_{\gLie}/(4\pi^2)^2$. Thus the candidate tensors are themselves
independent, and the kernel of divergence has dimension at most one.

The Maxwell stress is conserved and
\[
 -T^{\mathrm{Max}}
 =\tfrac12 Z^{(1)}-\tfrac14 Z^{(2)}+Z^{(3)}.
\]
It is nonzero by the preceding off-diagonal correlation. Hence it spans that
kernel.
\end{proof}

This result fixes only a line of conserved tensors. Multiplying the
stress by a scalar preserves all the properties used in the proposition.
We next fix its normalization by a local Ward flux and then classify
linear responses on fields of arbitrary bounded degree.


\section{Local stress fluxes}\label{ch03:ward}

The sign in \eqref{ch03:stress} is fixed by a local identity, not by
conservation. Let $x_0,\ldots,x_n$ be distinct points and choose a real
$\chi\in C_c^\infty(\RR^4)$ equal to one near $x_0$, with support
avoiding $x_1,\ldots,x_n$. For a vector field $v$ define the flux
\begin{equation}\label{ch03:flux}
 \mathcal F_{\chi,v}(P_0;P_1,\ldots,P_n)
 =-\sum_{a,b}\int \partial_a\chi(y)v_b(y)
 S_{n+2}^{\mathrm{Max}}(T^{\mathrm{Max}}_{ab},P_0,\ldots,P_n;
                       y,x_0,\ldots,x_n)\,dy.
\end{equation}
The integrand is supported on $\operatorname{supp}\nabla\chi$, a compact set separated from all the insertions.
It is therefore an ordinary integral of a smooth function; no extension
of a product through $x_0$ is involved.

For a constant translation $v=e_b$, let $\delta_v=\partial_b$ act on
local polynomial labels. For $v(y)=K(y-x_0)$, $K^{\mathsf T}=-K$,
let $\delta_v=\mathcal L_K$ be the infinitesimal action induced on jets
by the pullback of two-forms under $y\mapsto x_0+e^{tK}(y-x_0)$.
Equivalently, on a smooth closed two-form $f$ its variation is
\begin{equation}\label{ch03:lie}
 (\mathcal L_v f)_{ij}
 =v_b\partial_bf_{ij}+ (\partial_i v_b)f_{bj}
                            +(\partial_jv_b)f_{ib}
 =(d\iota_v f)_{ij},
\end{equation}
with sums over $b$. On polynomial labels the action is the induced
Leibniz derivation, including the derivative indices of each jet.
For the dilation $v(y)=y-x_0$, this action at the center is the weight
operator
\begin{equation}\label{ch03:weight}
 \mathscr E=\sum_{j,\lambda}(j+2)X_{j,\lambda}
                                  \frac{\partial}{\partial X_{j,\lambda}},
\end{equation}
where $X_{j,\lambda}$ are any linear coordinates on order-$j$ solution
jets, including the $\gLie$ components. The sum is finite on each polynomial.

\begin{proposition}[Normalized free Ward identities]\label{ch03:free-ward}
For translations, centered rotations and centered dilation,
\begin{equation}\label{ch03:ward-identity}
 \mathcal F_{\chi,v}(P_0;P_1,\ldots,P_n)
 =S_{n+1}^{\mathrm{Max}}(\delta_vP_0,P_1,\ldots,P_n;x_0,\ldots,x_n).
\end{equation}
This holds for arbitrary finite Wick polynomials and arbitrary finite
spectator lists. The common coefficient is $+1$ for the stress
\eqref{ch03:stress} and the outward-flux convention \eqref{ch03:flux}.
\end{proposition}

The Ward mechanism is the standard free-field action of the stress
current. We include the normalization and the treatment of two
contractions at the enclosed insertion, since neither is supplied by a
formal integration by parts through a singularity. Wick's theorem is
used throughout in its usual form~\cite{Wick1950}.

\begin{proof}
Polarize our stress between closed two-forms $f,g$:
\[
 t_{ab}(f,g)=-\sum_c\bigl(\langle f_{ac},g_{bc}\rangle_{\gLie}+\langle g_{ac},f_{bc}\rangle_{\gLie}\bigr)
             +\tfrac12\delta_{ab}\sum_{c,d}\langle f_{cd},g_{cd}\rangle_{\gLie}.
\]
The Bianchi identity gives
\[
 \sum_a\partial_at_{ab}
 =-\sum_c\big(\langle (\operatorname{div}f)_c,g_{bc}\rangle_{\gLie}
                         +\langle f_{bc},(\operatorname{div}g)_c\rangle_{\gLie}\big).
\]
The tensor is symmetric and traceless. Thus contracting with a constant,
skew-linear or scalar-linear $v$ produces no additional divergence term
from $Dv$.

Suppose $f$ is a smooth Maxwell wave near $\operatorname{supp}\chi$,
and $g$ is the covariance pulse from $F_{ij}(x_0)$. Regularize its Green
function by convolution with a smooth compactly supported unit-mass
approximate identity $\rho_\epsilon$, so
$-\Delta\Gamma_\epsilon=\rho_\epsilon$.
The resulting closed pulse has divergence
\[
 (\operatorname{div}g_\epsilon)_c
   =\delta_{cj}\partial_i\rho_\epsilon(\,\cdot-x_0)
       -\delta_{ci}\partial_j\rho_\epsilon(\,\cdot-x_0).
\]
Integrate the displayed divergence identity against $\chi v$.
An integration by parts against $\rho_\epsilon$ shows that the flux
converges to
$\partial_i(v_bf_{bj})(x_0)-\partial_j(v_bf_{bi})(x_0)
=(\mathcal L_vf)_{ij}(x_0)$.
Here $\chi=1$ near the center. On $\operatorname{supp}\nabla\chi$,
the regularized covariance $g_\epsilon$ and its derivatives of each fixed
order converge uniformly to $g$ and the corresponding derivatives as
$\epsilon\to0$. Differentiating in the pulse center, while
holding $v,\chi$ fixed, gives the same identity for every jet probe.
For dilation, \eqref{ch03:lie} becomes
$(y-x_0)\cdot\nabla f+2f$; its order-$j$ jet at $x_0$ is
$(j+2)$ times the original jet. This gives \eqref{ch03:weight}.

In a Wick correlation, extract the two factors at the stress insertion.
The remaining sum has three types. If both partners are outside $x_0$,
the two waves are smooth throughout the cutoff support and conservation
gives zero flux. If exactly one partner is at $x_0$, the preceding
calculation differentiates that occurrence. Summing its possible
partners reconstructs Wick's rule, and summing occurrences gives the
Leibniz derivation $\delta_v$.

If both partners are at $x_0$, their jet orders $j,l$ give pulse
homogeneities $-4-j$ and $-4-l$. Conservation away from $x_0$ makes
the flux independent of the choice of $\chi$. Rescaling a cutoff by $t$
about $x_0$ multiplies this same flux by $t^{e-5-j-l}$, where
$e=0$ for translation and $e=1$ for rotation or dilation.
Taking $t=2$ forces the flux to vanish. The flux integral involves the two
pulses only away from their common singularity at $x_0$.
The polarization counts each ordered pair twice, canceling its factor
$1/2$. Combining the three types proves \eqref{ch03:ward-identity}.
\end{proof}

\section{Separating fields by insertions of bounded scaling dimension}\label{ch03:bounded-detector}

The polynomial labels and their scaling dimensions are those of Section~\ref{ch02:jets-section}; throughout this section they are evaluated in the free curvature field defined in Section~\ref{ch03:covariance}.

Corollary~\ref{ch03:detectors} detects a field of degree $d$ by another
field of degree $d$. A different statement is needed when a charge can
raise degree: its response must be tested by spectators whose degree
remains fixed. Their \emph{number} may be arbitrary.

We first make explicit the polynomial space used here. The free image
$\mathcal M_d^{\mathrm{Max}}$ in \eqref{ch03:free-space} is the full
space of invariant polynomials of weight at most $d$ on the solution
jets. Indeed, the radial map identifies each order of solution jets with
a linear subspace of the corresponding radial-potential space.
Choose an invariant inner product and its orthogonal projection onto
this subspace, separately in each weight. Composing a polynomial with
these projections extends it to the radial jets without changing its
weight or invariance. This is the usual compact-group averaging
argument; no equality with a larger space of unrestricted formal jets
is being assumed.

Let $\mathfrak P$ be the polynomial algebra on all solution jets and
$\mathfrak P^G$ its invariant subalgebra. For $D\ge0$, define
\begin{equation}\label{ch03:ND}
 \mathcal N_D=\{P\in\mathfrak P:
 S^{\mathrm{Max}}(P(x_0),P_1(x_1),\ldots,P_n(x_n))=0\},
\end{equation}
where the equality is required for every separated word with
$P_i\in\mathcal M_D^{\mathrm{Max}}$. There is no restriction on
the degree of $P$ or the length $n$.

\begin{proposition}[A separating family of insertions of bounded scaling dimension]
\label{ch03:uniform-detector}
There is a finite $D_*$ such that
$\mathcal N_{D_*}\cap\mathfrak P^G=\{0\}$.
\end{proposition}

\begin{proof}
First restrict all jet orders to at most $j$. Hilbert's finite-generation
argument~\cite{Hilbert1890}, using compact-group averaging, gives
finitely many invariant generators homogeneous in ordinary polynomial
degree. Let $b_j$ bound their weights. To apply that standard theorem,
the jet space is finite-dimensional, the polynomial ring is Noetherian,
and averaging preserves its invariant ideal and homogeneous pieces.
Thus an invariant polynomial of ordinary degree $k$ is a sum of products
of generators of total ordinary degree $k$, each of weight at most $b_j$.

If a polynomial in $\mathcal N_{b_j}$ is nonzero, take its lowest
nonzero ordinary-degree part $P_k$. Express its reflected conjugate as
such a sum of products, and put the factors at distinct spectator points
approaching a common point different from $x_0$. There are precisely $k$
spectator legs. Higher-degree parts of $P$ cannot pair with them, and
for $P_k$ every spectator leg pairs with a leg at $x_0$. No singular
spectator--spectator contraction survives. The limit is the strictly
positive two-point pairing of $P_k$ with its reflected conjugate,
contradicting \eqref{ch03:ND}. This proves detection at bounded jet order.

To obtain a uniform bound, we use a collision that lowers the highest
jet order. For a linear functional $p$ on scalar order-$j$ jets, let
$X_j(p)\in\gLie$ be its evaluation on the spacetime indices of the jet.
Write $X_{j,\alpha}(p)=\langle e_\alpha,X_j(p)\rangle_{\gLie}$ in
the basis of Section~\ref{ch01:gauge}. Fix a polynomial
$P$ with highest jet order $j$ and a quadratic spectator
${:}\langle X_k(p),X_l(q)\rangle_{\gLie}{:}$, where $l<k$.
Contracting just its order-$k$ leg with the top-order variables of $P$
leaves the polynomial
\begin{equation}\label{ch03:replacement}
 \mathcal R_{p,q,j}(z)P
 =\sum_{\alpha,\lambda}X_{l,\alpha}(q)
 C_{k,j}(p,e_\lambda;z)
           \frac{\partial P}{\partial X_{j,\alpha}(e_\lambda)}.
\end{equation}
The covariance $C_{k,j}$ is in a single unit direction in $\gLie$. We claim that
\begin{equation}\label{ch03:replacement-null}
 P\in\mathcal N_D,\quad k+l+4\le D
 \quad\Longrightarrow\quad
 \mathcal R_{p,q,j}(z)P\in\mathcal N_D\quad(z\ne0).
\end{equation}

Here is the necessary separation of singularities. A differentiated
curvature covariance of orders $k,j$ is harmonic, homogeneous of degree
$-r$, $r=4+k+j$, and becomes smooth after multiplication by
$|z|^{2(r-1)}$. A nonzero such harmonic function cannot become smooth
after multiplication by $|z|^{2(r-2)}$. Indeed, for a harmonic
homogeneous $H$ of degree $-r$,
\[
 \Delta(|z|^{2s}H)=4s(s+1-r)|z|^{2s-2}H.
\]
Starting with a smooth extension at $s=r-2$ and iterating would extend
$H$ smoothly to zero, contradicting its negative homogeneity unless it
vanishes. A product of two harmonic poles of orders $r_1,r_2$,
on the other hand, is cleared by $|z|^{2(r_1+r_2-2)}$.

Insert the quadratic spectator at $x_0+tz$ and keep all other spectators
fixed. Its two contractions with $P$ give finitely many homogeneous
product poles. Subtract them before multiplying by $t^r$.
The limit of the remaining expression is the correlation with
\eqref{ch03:replacement}: the lower single-contraction orders vanish,
and contractions with remote spectators are smooth. The original full
correlation is zero by \eqref{ch03:ND}. Successively removing the
highest powers of $t^{-1}$ shows that the surviving harmonic pole is
minus a product-pole coefficient of the same degree, or is zero.
The extra radial cancellation just proved forces it to be zero in
either case. This proves \eqref{ch03:replacement-null}; in particular
we have not assumed existence of the unmodified rescaled collision limit.

Put $g=d_{\gLie}$ and choose
$D_*\ge\max(b_g,2g+5)$. Use $k=g+1$ and $l=0,\ldots,g-1$.
For each $l$ choose a nonzero probe $q_l$ and a solution jet on which it
has value one. Such jets exist in every order: for example, the curl of
$e_3\operatorname{Re}(x_1+ix_2)^{l+1}$ is a nonzero homogeneous
Maxwell solution of degree $l$. If the highest order $j$ of $P$ exceeds
$g$, each polynomial \eqref{ch03:replacement} has smaller weight than
$P$. Strong induction on weight and \eqref{ch03:replacement-null}
therefore make it zero.

The matrix of polynomial entries $(X_{l,\alpha}(q_l))_{l,\alpha}$
has nonzero determinant: choose the jets so that $X_l(q_l)=e_{l+1}$
for each $l$. At these jets the matrix is the identity. Multiplication
by its adjugate, in the polynomial integral domain, shows that every
componentwise top-order contraction derivative in
\eqref{ch03:replacement} vanishes. As $p,z$ vary, the functionals
$C_{k,j}(p,\cdot;z)$ span the dual of the order-$j$ solution space.
To verify this last assertion, a vector annihilated by all of them has
a response wave whose order-$k$ jets vanish everywhere off zero.
It is then a polynomial of degree below $k$ there; its negative
homogeneity makes it zero. Formula~\eqref{ch03:kelvin} makes the
original vector zero. Hence all top-order partial derivatives of $P$
vanish. Descending removes every jet order above $g$, and the fixed-order
argument with $b_g$ concludes the induction.
\end{proof}

\section{Finite correlation criteria for infinitesimal symmetries}\label{ch03:finite-ward}

We now give the two finite-dimensional consequences used later to
normalize translations and remove an internal rotation defect.
Throughout this section the Lie algebra is simple and nonabelian.
Recall that $\mathcal M_d^{\mathrm{Max}}$ consists of invariant
polynomials of weight at most $d$ on the solution jets $\mathcal S_j$;
its elements are local Wick fields, not distributions defined on collision
sets. For maps $A$ on this space, a \emph{Ward sum} means
\[
 W_A(P_1,\ldots,P_n;x)=\sum_{i=1}^n
 S_n^{\mathrm{Max}}(P_1,\ldots,AP_i,\ldots,P_n;x).
\]
All points in this definition are distinct. The word length is arbitrary.

\begin{theorem}[Characterization of infinitesimal translations]\label{ch03:vector-tests}
Let $d\ge\max(4,D_*)$. Consider real-preserving complex-linear maps
$A_b:\mathcal M_d^{\mathrm{Max}}\to\mathcal M_{d+1}^{\mathrm{Max}}$
which kill the unit, send $\mathcal M_e^{\mathrm{Max}}$ into
$\mathcal M_{e+1}^{\mathrm{Max}}$ for $e\le d$, and transform as a vector
under signed coordinate permutations. Their Ward sums vanish identically
if and only if $A_b=c\partial_b$ for a common real scalar $c$.
There are finitely many evaluations of configuration derivatives of these
Ward sums, at collinear configurations with line axis different from
$b$, whose common kernel is precisely this translation line.
\end{theorem}

\begin{proof}
Translation invariance proves one direction. We prove the converse,
starting at weight four. Let $\mathcal S_j^{\mathrm{sc}}$ be the
scalar-valued solution jets, so $\mathcal S_j=\gLie\otimes
\mathcal S_j^{\mathrm{sc}}$, and write
$U_j=(\mathcal S_j^{\mathrm{sc}})^*$ for the scalar order-$j$ probes and
$Q(p,q)=\langle X_j(p),X_k(q)\rangle_{\gLie}$ for the corresponding
quadratic polynomial. The invariant metric on a simple Lie
algebra is unique up to scalar. Hence $\mathcal M_4^{\mathrm{Max}}$
is spanned by the unit and $Q(p,q)$ with $p,q\in U_0$;
weight five adds exactly the mixed quadratics with $p\in U_0,q\in U_1$.
Inversion of all coordinates acts by $(-1)^e$ at weight $e$.
The filtration and vector parity therefore imply that
$A_bQ(p,p)$ is a mixed quadratic. It has a unique representation by a
linear map $D_p:\mathcal S_0^{\mathrm{sc}}\to U_1$:
\[
 A_bQ(p,p)=\sum_\lambda Q(e^\lambda,D_p(e_\lambda)).
\]
Uniqueness follows by assigning independent zero- and first-order
solution jets in a single unit direction in $\gLie$. Here and until the coefficient comparison,
$b$ is fixed.

Denote the scalar covariance of probes by $C(p,q;z)$.
The three-quadratic Wick rule is
\begin{align*}
&S_3^{\mathrm{Max}}(Q(a,v),Q(p,p),Q(q,q);x,y,z)\\
&\quad=4d_{\gLie}\bigl[C(a,p;x-y)C(v,q;x-z)
             +C(v,p;x-y)C(a,q;x-z)\bigr]C(p,q;y-z).
\end{align*}
In the Ward sum for three quadratics at $tu,0,z$, multiply by $t^5$
and let $t\downarrow0$. The third response has no fifth-order pole.
The one-jet covariance is odd, so the other two terms have opposite
signs. The zero-order covariance is invertible, allowing the probe at
$z$ to be chosen so its contraction against any $q$ is $q(X)$,
for any prescribed $X\in\mathcal S_0^{\mathrm{sc}}$. Consequently
\begin{equation}\label{ch03:panel}
 C(D_p(X),q;u)q(X)=C(D_q(X),p;u)p(X).
\end{equation}
If $p(X)=0$, varying $q$ with $q(X)\ne0$ shows that $D_p(X)=0$.
The separation used here follows from \eqref{ch03:kelvin}: a probe
pairing to zero with all zero-order probes at every displacement has
zero response wave, hence is zero. Thus
$D_p(X)=p(X)f_b(p)$ for some $f_b(p)\in U_1$.
Cancellation in \eqref{ch03:panel} gives
\[
 C(f_b(p),q;u)=C(f_b(q),p;u).
\]
This identity and separation prove additivity and real homogeneity of
$f_b$; the original signed-coordinate covariance passes to $f_b$.

For completeness the remaining coefficient determination is very small.
Write $c_{ij}$ for the probe $F_{ij}$ and $\partial_a c_{ij}$ for its
first derivative. Parity allows only
$\partial_1c_{12},\partial_3c_{23},\partial_4c_{24}$ in $f_1(c_{12})$.
The swap of axes $3,4$ and the divergence equation reduce their span to
$\partial_1c_{12}$. Similarly the candidates for $f_1(c_{23})$ are
$\partial_1c_{23},\partial_2c_{13},\partial_3c_{12}$;
the swap of axes $2,3$ and the Bianchi equation reduce them to
$\partial_1c_{23}$. The two coefficients agree by reciprocity:
differentiation of the covariance gives
$C(\partial_1c_{12},c_{23};e_1+e_3)=-1/(2\pi^2)\ne0$.
Signed permutations then give $f_b=c_0\partial_b$ on every probe.
Since $\partial_bQ(p,p)=2Q(p,\partial_bp)$, we have
$A_b=c\partial_b$ on weight four, with $c=c_0/2$.

Set $B_b=A_b-c\partial_b$. It kills the stress and still has zero
Ward sums. Write its Ward sum for a word with an additional stress
insertion. Apply the flux formula \eqref{ch03:flux} to this insertion
with the dilation vector field $v(y)=y-x_0$. The term acting on the stress is zero. The free Ward identity
then gives the Ward sum with $\mathscr E B_bP$ in the distinguished slot
and $\mathscr EP$ in the other terms. Subtract the ordinary Ward sum
for the word starting with $\mathscr EP$. It follows that
$[\mathscr E,B_b]P\in\mathcal N_d$. Proposition
\ref{ch03:uniform-detector} makes this commutator zero, even though its
weight can be $d+1$. Thus $B_b$ preserves each weight eigenspace.
But vector covariance makes it anticommute with coordinate inversion,
which is the scalar $(-1)^e$ on that eigenspace. Therefore $B_b=0$.

Finally, if every configuration derivative of a Ward sum vanishes at
all strict collinear configurations along one axis other than $b$,
real analyticity implies that it vanishes everywhere. One can connect
any separated configuration to such a collinear configuration without
collisions: order its projections on that axis, break ties by tiny
ordered perturbations, and interpolate to the resulting line points.
The projected differences keep a strict sign until the separated endpoint.
Analytic continuation along the segment proves the assertion.
The space of permitted maps is finite-dimensional over $\RR$.
An intersection of kernels of linear functionals on it is already the
intersection of a finite subfamily (take a basis of their span in the
dual). Real and imaginary parts of the evaluations give the required
finite family.
\end{proof}

\begin{theorem}[A vanishing criterion for an internal rotation defect]\label{ch03:internal-tests}
For $d\ge4$, let
$\Delta:\mathfrak{so}(4)\times\mathfrak{so}(4)
\to\operatorname{End}_{\CC}(\mathcal M_d^{\mathrm{Max}})$
be real-bilinear and alternating. Assume its values preserve reality and
the weight filtration and kill the unit, and that
\[
 \Delta(sKs^{-1},sLs^{-1})=\rho(s)\Delta(K,L)\rho(s)^{-1}
\]
for orientation-preserving signed coordinate permutations $s$.
If every value of $\Delta$ has zero Ward sums, then $\Delta=0$.
It suffices to check that finitely many of these Ward sums vanish.
\end{theorem}

\begin{proof}
Fix $K,L$ and put $U=\Delta(K,L)$. At weight four write
$B_x(P,Q)=S_2^{\mathrm{Max}}(P,Q;x,0)$.
Recall $R_x=I-2xx^{\mathsf T}/|x|^2$, the reflection normal to $x$.
The zero-order covariance is
$\pi^{-2}|x|^{-4}\bigwedge^2R_x$ tensored with the identity on $\gLie$.
Wick's rule consequently gives, for $|x|=|y|$,
\[
 B_x(P,Q)=B_y(\rho(R_yR_x)P,Q)
 \qquad(P,Q\in\mathcal M_4^{\mathrm{Max}}).
\]
The two-point Ward sum says that $U$ is skew for every $B_x$.
These forms are nondegenerate by strict reflected positivity.
Comparing skewness for $x$ and $y$ therefore shows that $U$ commutes
with $\rho(R_xR_y)$ on this space. In particular it commutes with each
coordinate quarter-turn $q_{ab}$, by choosing $x$ parallel to $e_a$
and $y$ parallel to $e_a+e_b$.

An alternating bilinear map on $\mathfrak{so}(4)$ invariant under
simultaneous conjugation by these quarter-turns is zero, even when its
values lie in an arbitrary real vector space. To see this on the basis
$E_{ab}=e_ae_b^{\mathsf T}-e_be_a^{\mathsf T}$, a half-turn in the
plane spanned by the third and fourth axes fixes one member of a pair
sharing an axis and negates the other. For a disjoint pair, the even
permutation exchanging their two ordered axis pairs exchanges the
arguments, so alternation gives zero. Equal pairs vanish by alternation.
Bilinearity proves the claim. Covariance and the commutation just proved
apply this observation to $\Delta(K,L)P$ for every weight-four $P$.
In particular $U$ kills the stress.

For each $H\in\mathfrak{so}(4)$, repeat the preceding commutator
argument using the rotation vector field $v(y)=H(y-x_0)$ in the flux
formula \eqref{ch03:flux}. This gives $[\mathcal L_H,U]P\in\mathcal N_d$.
This time the commutator itself lies in $\mathcal M_d^{\mathrm{Max}}$,
so a single reflected two-point pairing shows that it vanishes. The uniform
bound $D_*$ is unnecessary. Exponentiation in the finite-dimensional
space shows that $U$ commutes with every quarter-turn on the whole
space. Covariance and the same alternating-map argument give
$\Delta(K,L)P=0$ for every $P$. The finiteness assertion again follows
by taking a basis of the span of the Ward functionals in the dual of
the finite-dimensional space of allowed $\Delta$.
\end{proof}

Each criterion involves finitely many evaluations on a specified space of possible responses. Their
application to the interacting construction requires proving that its
responses obey the stated filtration, covariance, and comparison
hypotheses. No such interacting input is part of the free-field results
proved here.
